%% Effective Training Pressure Gates Recursive Knowledge Degradation in LLMs
%% Target: Engineering Applications of Artificial Intelligence (Elsevier)
%% Template: CAS Double Column (cas-dc)

\documentclass[a4paper,fleqn]{cas-dc}
\usepackage[authoryear,longnamesfirst]{natbib}
\usepackage{booktabs}
\usepackage{amsmath}

\begin{document}
\let\WriteBookmarks\relax
\def\floatpagepagefraction{1}
\def\textpagefraction{.001}

\shorttitle{Pressure-Gated Recursive Degradation in LLMs}
\shortauthors{Azancort et~al.}

\title[mode=title]{Effective Training Pressure Gates Recursive Knowledge
  Degradation in LLMs: A Multi-Axis Dose-Response Study}

\author[1]{Julio Azancort}[
  orcid=0000-0001-7511-2910,
  role=Researcher]
\cormark[1]
\ead{jazancort@ufpa.br}
\credit{Conceptualization, Methodology, Software, Investigation,
  Writing - Original Draft, Visualization}

\affiliation[1]{organization={Universidade Federal do Par\'a (UFPA)},
    city={Bel\'em},
    state={PA},
    country={Brazil}}

\cortext[cor1]{Corresponding author}

% ============================================================
% ABSTRACT
% ============================================================
\begin{abstract}
Recursive training on synthetic data causes model collapse, but the transition
from stability to degradation is sharper and more controllable than previously
understood. We show that recursive factual degradation under QLoRA is governed
by effective training pressure, a composite of adapter capacity, update
magnitude, and synthetic exposure. Through systematic dose-response mapping
across three backbones (Qwen 2.5 1.5B, Gemma 3 1B, Gemma 4 E2B), full
fine-tuning comparison, and causal intervention experiments, we find: (1) a
sharp regime transition between homeostatic retention and progressive
degradation; (2) backbone-dependent thresholds not reducible to a universal
effective-rank value; (3) a one-time adaptation cost advantage of approximately
5 percentage points for QLoRA over full fine-tuning at matched perturbation; (4)
output drift as a diagnostic signature of above-threshold regimes; and (5)
restoration of near-homeostatic stability by reducing synthetic exposure by
approximately 5\%, confirmed across three independent random masks. These
results position recursive degradation as pressure-gated, with practical
implications for monitoring and governing recursive training pipelines.
\end{abstract}

% ============================================================
% HIGHLIGHTS
% ============================================================
\begin{highlights}
\item Recursive degradation is governed by effective training pressure, not adapter rank alone
\item A sharp backbone-dependent boundary separates homeostatic and degradative regimes
\item Reducing synthetic exposure by approximately 5\% restores near-homeostatic stability across three independent random masks
\item Low-rank QLoRA reduces one-time adaptation cost versus full fine-tuning by approximately 5 percentage points
\item Output drift (verbosity, efficiency loss) serves as a diagnostic signature of above-threshold regimes
\end{highlights}

% ============================================================
% KEYWORDS
% ============================================================
\begin{keywords}
recursive synthetic training \sep model collapse \sep
parameter-efficient fine-tuning \sep LoRA \sep
effective training pressure \sep knowledge retention
\end{keywords}

\maketitle

% ============================================================
% 1. INTRODUCTION
% ============================================================
\section{Introduction}\label{sec:introduction}

% TODO: Write introduction

% ============================================================
% 2. RELATED WORK
% ============================================================
\section{Related Work}\label{sec:related}

% TODO: Write related work

% ============================================================
% 3. METHODOLOGY
% ============================================================
\section{Methodology}\label{sec:methodology}

We study the stability of recursive synthetic fine-tuning through the lens of
\textit{effective training pressure}: the combined effect of update capacity,
perturbation magnitude, and synthetic-data exposure applied at each recursive
generation. We do not treat this quantity as a single fitted scalar. Instead, we
use it as an organizing framework in which each component is manipulated
separately: adapter rank controls update capacity, learning rate controls
perturbation magnitude, and filtering or downsampling controls synthetic
exposure. This design lets us ask whether transitions between homeostatic and
degradative regimes are attributable to capacity, magnitude, exposure, or their
interaction.

Our experiments follow a dose-response design across these three axes, tracking
factual retention and output-distribution drift over up to ten recursive
generations. In the main protocol, each generation is trained only on synthetic
outputs produced by the previous generation, with no real-data accumulation,
replay buffer, or quality filter. This strict replace-only setting follows the
recursive data-replacement paradigm studied in the model-collapse
literature~\citep{shumailov2024,dohmatob2025,gerstgrasser2024}, while allowing
us to isolate how the update mechanism changes the resulting dynamics.

The design is motivated by both practical and theoretical considerations. In
practice, parameter-efficient fine-tuning via LoRA~\citep{hu2022} and
QLoRA~\citep{dettmers2023} dominates production adaptation of large language
models, yet behavior under recursive synthetic exposure remains poorly
characterized beyond the observation that LoRA can ``learn less and forget
less''~\citep{biderman2024}. That work does not map when this regularization
remains protective and when it fails under repeated synthetic training.
Theoretically, fine-tuning updates often lie in low-dimensional
subspaces~\citep{aghajanyan2021}, suggesting that adapter rank directly
modulates the effective dimensionality of the perturbation applied each
generation. We operationalize this via the effective rank of the composed update
matrices $BA$, following the entropy-based definition of
\citet{roy2007}.

The remainder of this section is organized as follows: the recursive protocol
(\S\ref{sec:protocol}), factual retention benchmark (\S\ref{sec:k0}), model
backbones (\S\ref{sec:models}), experimental axes (\S\ref{sec:axes}),
effective-rank and perturbation metrics (\S\ref{sec:pressure}),
output-distribution diagnostics (\S\ref{sec:metrics}), synthetic-exposure
interventions (\S\ref{sec:intervention}), and statistical conventions
(\S\ref{sec:statistics}).

\subsection{Recursive replace-only protocol}\label{sec:protocol}

% TODO: Write protocol subsection

\subsection{K0 factual retention metric}\label{sec:k0}

% TODO: Write K0 metric subsection

\subsection{Models and hardware}\label{sec:models}

% TODO: Write models subsection

\subsection{Experimental axes}\label{sec:axes}

% TODO: Write axes subsection

\subsection{Effective training pressure}\label{sec:pressure}

% TODO: Write pressure concept subsection

\subsection{Output-distribution diagnostics}\label{sec:metrics}

% TODO: Write metrics subsection

\subsection{Synthetic-exposure interventions}\label{sec:intervention}

% TODO: Write intervention subsection

\subsection{Seeds, statistics, and reporting}\label{sec:statistics}

% TODO: Write statistics subsection

% ============================================================
% 4. RESULTS
% ============================================================
\section{Results}\label{sec:results}

% TODO: Write results

% ============================================================
% 5. DISCUSSION
% ============================================================
\section{Discussion}\label{sec:discussion}

% TODO: Write discussion

% ============================================================
% 6. LIMITATIONS
% ============================================================
\section{Limitations}\label{sec:limitations}

% TODO: Write limitations

% ============================================================
% 7. CONCLUSION
% ============================================================
\section{Conclusion}\label{sec:conclusion}

% TODO: Write conclusion

% ============================================================
% CREDIT
% ============================================================
\printcredits

% ============================================================
% REFERENCES
% ============================================================
\bibliographystyle{cas-model2-names}
\bibliography{cas-refs}

\end{document}
